\subsection{The sharp conductor-forty boundary}
\begin{lemma}[A root-of-unity length certificate]\label{lem:rootlength40}
Suppose that $\xi_1,\ldots,\xi_n$ are roots of unity and
\[
 \sum_{j=1}^n\xi_j=a+b\sqrt{10},\qquad a,b\in\mathbb Q.
\]
Then $n\geq -a-8b$. If equality holds, every $\xi_j$ is either $-1$ or
is a primitive fortieth root $\zeta_{40}^{r}$ with
$\chi_{40}(r)=-1$, where $\chi_{40}(r)=(40/r)$ is the primitive quadratic
Dirichlet character of conductor $40$.
Consequently, if
\[
 n=18,\qquad\sum_{j=1}^{18}\xi_j=-2-2\sqrt{10},
\]
then exactly two of the roots are $-1$, and the normalized absolute trace of
$\sum_j\xi_j^2$ is $2$.
\end{lemma}

\begin{proof}
Put $s=\sqrt{10}$. Let $K=\mathbb Q(s)$, whose conductor is $40$, and choose a cyclotomic field
$L$ containing $K$ and all the roots under consideration. Write
\[
 T(u)=\frac{\operatorname{Tr}_{L/K}(u)}{[L:K]},\qquad
 \ell(a+bs)=-a-8b.
\]
We claim that $\ell(T(\xi))\leq1$ for every root of unity $\xi$, with precisely
the equality cases stated above. The normalized relative trace does not depend
on the choice of common cyclotomic overfield.

Let $m$ be the exact order of $\xi$. If $40\nmid m$, the intersection
$K\cap\mathbb Q(\xi)$ is $\mathbb Q$: otherwise the conductor of $K$ would
divide $m$. Restriction of automorphisms consequently gives
\[
 T(\xi)=\frac{\mu(m)}{\varphi(m)}.
\]
It follows that $\ell(T(\xi))\leq1$, with equality only when $m=2$.

Now suppose $m=40h$, and write $\xi=\zeta_m^r$, $(r,m)=1$. Since $\mu(m)=0$,
the rational coefficient of $T(\xi)$ is zero. The quadratic coefficient is
computed by the induced quadratic Gauss sum:
\[
 G_m(r)=\sum_{v\in(\mathbb Z/m\mathbb Z)^\times}
             \chi_{40}(v)\zeta_m^{rv}.
\]
For completeness, the standard induced-Gauss-sum identity gives
\[
 G_m(r)=
 \begin{cases}
 \chi_{40}(r)\mu(h)\chi_{40}(h)\,2s,&(h,40)=1,\\
 0,&(h,40)>1.
 \end{cases}
\]
Here $\mu(h)=0$ already covers nonsquarefree $h$. This identity can also be
obtained directly from the Chinese remainder theorem when $(h,40)=1$:
the sum factors into the primitive Gauss sum
$G_{40}(1)=\sqrt{40}=2s$ and the Ramanujan sum $\mu(h)$, with the displayed
quadratic character factors. If an extra power of $2$ or $5$ occurs in $m$,
summing over the last such prime-power layer gives zero.
The terms with $\chi_{40}(v)=1$ are the automorphisms fixing $K$; their sum
is $(\mu(m)+G_m(r))/2$. Since there are $\varphi(m)/2$ such automorphisms,
\[
 T(\xi)=\frac{G_m(r)}{\varphi(m)}
 =\frac{\chi_{40}(r)\mu(h)\chi_{40}(h)}{8\varphi(h)}s
 \quad\text{if }(h,40)=1,
\]
and it is zero otherwise. Thus $|\ell(T(\xi))|\leq1$. Equality with value
$1$ forces $\varphi(h)=1$ and $(h,40)=1$, hence $h=1$, and then
$\chi_{40}(r)=-1$. This proves the claim. Applying $\ell\circ T$ to the sum
of the roots proves the inequality and its equality statement.

In the particular equality case, primitive fortieth roots have normalized
absolute trace $\mu(40)/\varphi(40)=0$, whereas $-1$ has trace $-1$.
The absolute trace of $-2-2s$ is $-2$, so exactly two roots are $-1$.
The square of a primitive fortieth root is primitive of order $20$ and has
normalized absolute trace $\mu(20)/\varphi(20)=0$. The two other squares
are $1$, proving the last assertion.
\end{proof}
